\subsection*{问题 1}
\paragraph{(a)}\label{sol:gae-a}
$n$ 步优势等于前 $n$ 个折扣 TD 残差之和。有限混合必须将权重 $\lambda^{m-1}$ 赋给最长目标。若只保留 $(1-\lambda)\lambda^{n-1}$、$n\leq m$，权重和为 $1-\lambda^m$，仍有 $\lambda^m$ 未分配。当 $m=1$ 时，唯一目标权重为一。当 $\lambda=1$ 时，全部权重集中于 $m$ 步目标。

\paragraph{(b)}\label{sol:gae-b}
$\lambda=0$ 时，首步优势为 $\delta_0=0.86$。$\lambda=1$ 时，完整折扣回报为 $1+0.9\times2+0.9^2\times3=5.23$，减去 $V_0=0.5$ 得到 $4.73$。两者分别对应一步自举目标和本条终止轨迹的完整回报目标。

\paragraph{(c)}\label{sol:gae-c}
时刻 $k$ 的残差误差为 $\gamma\Delta$。此前每向前一项，误差乘以 $\gamma\lambda$，所以第 $k-j$ 项误差为 $\gamma(\gamma\lambda)^j\Delta$。遇到先前的轨迹边界，传播由 $c=0$ 中断。较小的 $\lambda$ 减小这类远端误差的权重；这一结论只讨论固定价值误差的传播，不保证整体优势估计方差或训练回报单调改善。

\paragraph{(d)}\label{sol:gae-d}
从缓冲区末端开始令尾项为零。每步先根据真正终止标记计算自举，再根据两类标记控制递推。返回 $A_k$ 与 $A_k+V_k$。测试需使用数值已知的例题和独立混合表达式。时间截断与真正终止的输入可具有相同形状，单靠形状检查无法区分它们的语义。

\paragraph{(e)}\label{sol:gae-e}
对每个起点单独计算全部 $n$ 步目标，再按有限混合权重组合。将其与后向实现逐项比较，避免两种实现共用同一个递推函数。价值替换实验应符合 (c) 的解析误差。此测试覆盖有限混合、边界与误差传播；完整策略训练的稳定性和回报需要另外运行环境与更新循环。
